\chapter{Generating Type Systems: OSLF and the Hypercube}
\label{ch:hypercube}

\section{OSLF as a family}

Operational Semantics in Logical Form (OSLF) is not a single algorithm but a
\emph{family} of them. Each member takes a theory presented as $(\Sigma, E, R)$ together
with a specification of a target logic, and generates from the two an assignment of
predicates to terms. The lineage runs through Hennessy--Milner
logic~\cite{HennessyMilner1985}, Stirling's modal characterisations, and, for the
reflective setting, the namespace logic of the rho calculus~\cite{NamespaceLogic}.

The member we develop here is the \emph{Hypercube} algorithm of Stay, Meredith and
Wells, which generates not a modal logic for observing programs but a family of type
systems for classifying them. Both directions matter, and it is worth saying why they are
the same construction. A type is a predicate on terms; a modal formula is a predicate on
terms; the difference is bookkeeping about what one intends to do with the predicate. What
the family provides is the machinery for reading a rewrite rule as a formation rule for
predicates.

\section{Adequacy is not automatic}
\label{ob:subsec:adequacy}

Because the difference between the two directions is bookkeeping, it is tempting to
suppose that anything the family generates inherits the property one wants from a
generated modal logic --- \emph{adequacy}, the characterisation of the theory's
behavioural equivalence by satisfaction. It does not, and the reason is worth stating
before the construction rather than after it.

Adequacy is a property of the \emph{target logic specification}, not of the generating
machinery. The classical Hennessy--Milner argument needs a certain minimum of logical
structure to run: something akin to a conjunction, so that a formula can record several
observations of the same state at once, and something akin to a negation, so that a
formula can record the \emph{absence} of an observation. Strip either and the argument
fails --- without negation one characterises a simulation preorder rather than
bisimulation; without conjunction one cannot separate states that differ only in how
their capabilities are bundled. A generated logic lacking either connective will fail to
have an adequacy theorem, and no amount of care in the generator repairs this, because
the deficiency is in what one asked the generator to produce.

\begin{remark}[The Hypercube functor offers no adequacy theorem]
\label{ob:rem:no-adequacy}
The Hypercube algorithm generates a \emph{type system}, and type systems are not
generally expected to be adequate. Nobody asks of the simply-typed $\lambda$-calculus
that its typing judgement characterise observational equivalence; one asks that it be
sound, decidable, and useful. The same standard applies here. Sections~\ref{ob:subsec:modal}
onwards should therefore be read as generating a classification discipline, not as
generating a behavioural characterisation.
\end{remark}

This is a deliberate trade, and Section~\ref{ob:subsec:dials} describes what is bought with
it: the user dials which connectives are admitted, precisely in order to expose a surface
on which a finitely checkable fragment is generable. The cost of that expressiveness and
control is that not everything generated is necessarily adequate.

\begin{remark}[And adequate for what?]
\label{ob:rem:adequate-for-what}
There is a second question hiding in the word, which the discussion above does not
touch. Even granting all the connectives the Hennessy--Milner argument needs, the
generated logic has a \emph{structural} layer that the argument says nothing about,
because a decomposition of a term is not a step it can take. That layer is in
general strictly more discriminating than the behavioural equivalence one would
naturally measure it against --- which, taken at face value, makes adequacy
impossible rather than merely unguaranteed.

It is not impossible; the two claims are about two different relations. But the
resolution needs a construction, and the construction is the subject of
Chapter~\ref{ch:obsext}. Until then, read every claim of adequacy in this book as
carrying an index that has not yet been written down.
\end{remark}

\section{Input: rewrites and redex positions}

The Hypercube algorithm takes as input a theory in the classifying form of
Section~\ref{ob:subsec:classifying}. Its base rewrites have the schematic shape
\[
  x_1{:}X_1, \dots, x_n{:}X_n \ \ctxbar\ \emptyset \ \vdash\ L(\vec x) \rsq R(\vec x) ,
\]
and --- this is the crucial extra datum --- it also tracks \emph{redex positions}.
Choosing a subterm occurrence $t_j$ of $L$ with carrier $Y_j$ determines a one-hole
context $K_j[-]$ with $K_j[t_j] = L$. Write
\[
  V_j = \fv(K_j) - \fv(t_j), \qquad W_j = \fv(t_j) - \fv(K_j) .
\]
The variables in $V_j$ --- free in the context but not in the chosen subterm --- become
the \emph{rely parameters} of the modality generated at that position.

\begin{remark}[The relationship to the interaction cut]
The reader will recognise $K_j$. In a theory presented as an interaction cut, the contact
site $K$ of $(\dagger)$ \emph{is} a one-hole context in exactly this sense, and the
distinguished positions are the two introductions. The interaction cut is thus the
special case in which one privileges a particular family of redex positions, and the
Hypercube algorithm is what one gets by not privileging them. This is why
Chapter~\ref{ch:igslt} does real work for the present section: it identifies which
positions are canonical.
\end{remark}

\section{The modal layer}
\label{ob:subsec:modal}

For each base rewrite and each chosen redex position, the algorithm introduces a primitive
modality
\[
  \angles{K_j}{\vec x :: \vec A}{B} ,
\]
a type former at carrier $Y_j$, read as a \emph{rely-possibly} specification: under the
rely assumptions $\vec x :: \vec A$, placing a term into the context $K_j[-]$ takes one
step to a specified right-hand side inhabiting $B$.

\begin{mathpar}
\inferrule*[right=M-Form]
  { \bigl(\Gamma \ctxbar \Delta \vdash A_k :: s_k^{X_k}\bigr)_{x_k \in V_j} \\
    \Gamma, \vec x{:}\vec X \ctxbar \Delta, \vec x{::}\vec A \vdash B :: s_{\mathrm{out}}^{\Pr} }
  { \Gamma \ctxbar \Delta \vdash \angles{K_j}{\vec x :: \vec A}{B} :: s_{\mathrm{out}}^{Y_j} }
\end{mathpar}

\begin{mathpar}
\inferrule*[right=M-Intro]
  { \bigl(\Gamma \ctxbar \Delta \vdash A_k :: s_k^{X_k}\bigr)_{x_k \in V_j} \\
    \Gamma, \vec w{:}\vec W, \vec x{:}\vec X \ctxbar \Delta, \vec x{::}\vec A \vdash B :: s_{\mathrm{out}}^{\Pr} \\
    \Gamma, \vec w{:}\vec W, \vec x{:}\vec X \ctxbar \Delta, \vec x{::}\vec A \vdash R(\vec x, \vec w) :: B }
  { \Gamma, \vec w{:}\vec W \ctxbar \Delta \vdash t_j(\vec w) :: \angles{K_j}{\vec x :: \vec A}{B} }
\end{mathpar}

Elimination is where a rewrite-generated modality differs from an equation-generated one.
Because the modality comes from a \emph{rewrite}, elimination does not provide full
conversion; it provides the operational step to the specified right-hand side, and the
typing of that right-hand side.

\begin{mathpar}
\inferrule*[right=M-Step]
  { \Gamma \ctxbar \Delta \vdash t :: \angles{K_j}{\vec x :: \vec A}{B} \\
    \bigl(\Gamma \ctxbar \Delta \vdash u_k :: A_k\bigr)_{x_k \in V_j} }
  { \Gamma \ctxbar \Delta \vdash K_j[t][\vec u/\vec x] \rsq R(\vec u) }
\and
\inferrule*[right=M-Elim]
  { \Gamma \ctxbar \Delta \vdash t :: \angles{K_j}{\vec x :: \vec A}{B} \\
    \bigl(\Gamma \ctxbar \Delta \vdash u_k :: A_k\bigr)_{x_k \in V_j} }
  { \Gamma \ctxbar \Delta \vdash R(\vec u) :: B[\vec u/\vec x] }
\end{mathpar}

To recover a conversion-like principle stable under rewriting, the algorithm freely adds
an \emph{empty-context possibility} operator $\Diamond$ at the carrier of programs, with
$p :: \Diamond B$ asserting that $p$ can take one step to some reduct of type $B$.
Equation-generated modalities, by contrast, do support full conversion by the usual rule ---
they are ``zero-step'' steps.

\section{The structural and propositional layers}

The modal fragment is intentionally behavioural, and this has a consequence that should
be stated plainly rather than discovered: \emph{terms with no available one-step behaviour
need not have any modal type at all.} A value, a normal form, a deadlocked process ---
none of these is classified by a layer whose formation rules are indexed by rewrites.

Two further layers repair this and extend expressivity.

\paragraph{Structural types.} For each term former $f$ the algorithm freely adds a type
former $f^\sharp$ whose inhabitants are exactly those terms whose root former is $f$.
These classify by the root of the parse tree, giving a baseline typing for programs which
can then be combined with the modal layer to express shape and one-step behaviour in the
same language. When the interaction constructor is read this way it becomes a
separating-style connective: a term satisfies $K(\varphi_1,\varphi_2)$ when it splits
across $K$ into a part satisfying $\varphi_1$ and a part satisfying $\varphi_2$. When $K$
is associative--commutative this is a genuine separating conjunction in the sense of
spatial and separation logics~\cite{CairesCardelli,Reynolds2002}; when $K$ is free it
is a positional split.

\paragraph{Propositional operations, stratified.} At each carrier and sort the algorithm
adds propositional operations on type formers --- but \emph{stratified by the categorical
strength required to interpret them}. The finite-limits fragment supplies $\top$ and
$\wedge$. Additional structure is required for $\bot$, for finite joins and
distributivity, and finally for implication and negation.

\begin{remark}[What a specification must supply]
\label{ob:rem:stratification}
The stratification is not a technicality of the metatheory; it is visible to the user of
a generated logic. A specification whose classifying theory has finite limits and no more
yields a conjunctive fragment; implication and negation become available only when the
theory supplies the structure to interpret them. A developer writing a property is
therefore working in a logic whose connectives are determined by their own specification,
and a property that fails to elaborate may be failing for this reason rather than because
it is false.
\end{remark}

\section{Sort slots and the equational centre}

Why a \emph{hypercube}? Each generated type former carries finitely many sort slots: one
for each rely input and one for the output. Filling each slot with $\ast$ or $\Box$ yields
a raw Boolean cube of choices. But equational axioms and rewrite laws force certain slots
to agree, and only the assignments stable under those laws survive. These form the
\emph{equational centre} $\Zc$, and the vertices of the resulting hypercube are the
elements of $\Zc$.

The cube's axes are not postulated. They are determined locally by the free variables of
operational contexts and the result types demanded by the rules.

\section{Worked example: a rho-calculus modality}
\label{ob:subsec:rho-modality}

Fix arity $1$. The communication rewrite is
\[
  n{:}\Nm,\ p{:}\Pr,\ \lambda x.q : [\Nm \to \Pr] \ \ctxbar\ \emptyset \ \vdash\
  \mathsf{out}_1(n,p) \mid \mathsf{in}_1(n, \lambda x.q) \ \rsq\ q[@p/x] .
\]
Choose the continuation redex position
\[
  t_j := \lambda x.q, \qquad K_j([-]) := \mathsf{out}_1(n,p) \mid \mathsf{in}_1(n,[-]) ,
\]
so that $K_j[t_j]$ is the left-hand side. The free variables of $K_j$ are
$V_j = \{n,p\}$, hence the induced modality is $\angles{K_j}{n :: A_n,\ p :: A_p}{B}$ at
carrier $[\Nm \to \Pr]$, contributing three local sort slots
\[
  \{\, s_n^{\Nm},\ s_p^{\Pr},\ s_{\mathrm{out}} \,\} .
\]
Relative to the origin assignment $(\ast,\ast,\ast)$, every nonempty subset of the three
slots can be lifted to $\Box$, yielding $2^3 - 1 = 7$ local axes.

Now hold $s_n = \ast$ fixed. Operationally the channel $n$ serves only to ensure that the
input and output are co-addressed; the reduct $q[@p/x]$ does not otherwise depend on it.
The remaining two-slot subgeometry is:

\begin{center}
\begin{tabular}{@{}lcc@{}}
\toprule
\textbf{Vertex (at $s_n = \ast$)} & $s_p^{\Pr}$ & $s_{\mathrm{out}}$ \\
\midrule
simply-typed corner        & $\ast$ & $\ast$ \\
polymorphism direction     & $\Box$ & $\ast$ \\
dependent-types direction  & $\ast$ & $\Box$ \\
type-constructors direction& $\Box$ & $\Box$ \\
\bottomrule
\end{tabular}
\end{center}

This is Barendregt's square~\cite{Barendregt1991}, recovered exactly. The three nontrivial
directions correspond to the three extra edges of the Lambda Cube beyond the simply-typed
corner. \emph{The Cube is one face of one modality of one rewrite rule.} It arises here
not because it was postulated but because the communication rule has two rely parameters
and one output, and the sort slots of a two-parameter modality form a square.

\begin{remark}[Why Barendregt needed no equational centre]
In the Lambda Cube there are no built-in equations, and the underlying equations of
$\beta$ add no relations between sorts. The whole raw Boolean cube is therefore
admissible, $\Zc$ is everything, and the distinction between the raw cube and the centre
never arises. It arises for us as soon as a specification carries equations --- which, by
rung two of the ladder, is essentially always.
\end{remark}

\section{The construction is a monad}

Let $\Lcat$ denote the category of theories in classifying form, with morphisms preserving
$\Pr$, $\rsq$, and the internal-language structure, and let $\Lcat_\Sigma$ denote the
category of such theories equipped with the additional structure provided by the modal,
structural, and propositional layers. On a morphism $F$, the action is by transport:
\[
  \angles{K_j}{\vec x :: \vec A}{B} \ \longmapsto\
  \angles{F(K_j)}{\vec x :: F(\vec A)}{F \cdot B} ,
  \qquad \Diamond B \longmapsto \Diamond (F \cdot B) .
\]

\begin{proposition}
There is a forgetful functor $U : \Lcat_\Sigma \to \Lcat$ and a left adjoint
$F : \Lcat \to \Lcat_\Sigma$ freely adjoining the layers, and the induced endofunctor
\[
  \hypK \;=\; U \circ F \;:\; \Lcat \longrightarrow \Lcat
\]
is a monad on $\Lcat$: for each theory $\mathsf{T}$, $\hypK(\mathsf{T})$ is the underlying
theory of its free extension by generated type formers.
\end{proposition}

This completes the pattern of Part~III. Three constructions, three monads, three
free/forgetful resolutions:
\[
  \Cost \ \text{(meter it)}, \qquad
  \Hist \ \text{(log it)}, \qquad
  \hypK  \ \text{(type it)} .
\]
In each case the left adjoint installs apparatus generated from the theory's own
presentation, and the right adjoint takes it away.

\section{Three dials}
\label{ob:subsec:dials}

Section~\ref{ob:subsec:adequacy} said that adequacy is a property of the target logic
specification. This subsection is the other half of that statement: the target logic
specification is something the \emph{user} supplies, and supplying it is the design
activity around which the whole apparatus is arranged.

Because the logic is generated from a language specification together with a choice of
connectives, a user can select \emph{fragments} to dial in finitely checkable properties.
There are three axes:

\begin{enumerate}[label=(\arabic*),leftmargin=2.5em]
  \item \textbf{The language fragment.} Restrict the sub-signature over which properties
    are stated. A property quantifying over a decidable fragment of the terms is a
    different proposition from the same property over the whole theory.
  \item \textbf{The connective fragment.} Restrict which of the generated connectives are
    in play, subject to Remark~\ref{ob:rem:stratification}. The conjunctive--modal fragment
    behaves very differently from one closed under negation and implication.
  \item \textbf{The cube vertex.} Choose an element of the equational centre $\Zc$. This
    is the classical design choice --- how much dependency to admit --- now presented as a
    coordinate rather than a fork in the literature.
\end{enumerate}

Each axis trades expressivity for checkability, and the three are independent. In the
general case, where the user has chosen not to constrain things statically, an
\emph{inference token budget} prevents runaway property checks. The budget is the backstop,
not the primary mechanism: the primary mechanism is fragment selection, and the budget
catches what selection did not.

\begin{remark}[Adequacy and checkability pull against each other]
\label{ob:rem:tension}
It would be convenient if the dials traded a property nobody wanted for one everybody did.
They do not, and the tension is structural rather than incidental. Adequacy requires
negation (Section~\ref{ob:subsec:adequacy}); negation is among the first connectives a user
drops when chasing decidability, since it is what turns a search for a witness into a
search over all possible witnesses. A user who dials down to a finitely checkable fragment
has, in the typical case, dialled out of the adequate region on the way.

The right conclusion is not that one of the two properties is illusory but that they
answer different questions. Adequacy asks whether the logic can \emph{distinguish} exactly
the programs the theory distinguishes --- a question about the logic's resolving power as a
whole. Checkability asks whether a particular property can be decided of a particular
program in bounded time. A user guarding a communication with a \texttt{where} clause,
in the manner of Chapter~\ref{ch:carho}, wants the second and has no use for the first. A user reasoning about when two
implementations may be substituted wants the first and will accept a semi-decision
procedure for it. The dials exist so that neither user has to pay for the other's
requirement.
\end{remark}

